\documentclass[12pt]{article}
\usepackage[margin=1in]{geometry}
\usepackage{amsmath}
\usepackage{setspace}
\usepackage{indentfirst}

\title{Optimal Fiscal and Monetary Policy in an Economy Without Capital\\
\large Section 6: Conclusions}
\author{Robert E. Lucas, Jr. and Nancy L. Stokey\\
\small Reproduced from Discussion Paper No. 532, August 1982}
\date{}

\begin{document}

\maketitle

\section*{6. Conclusions}

The main lessons to be drawn from the examples we have worked through in this paper were anticipated with remarkable accuracy and clarity by Alexander Hamilton in his \textit{Second Report on Public Credit} [1795]. The concluding pages of this report are devoted to arguing that the interest on government debt should not be taxable, and his arguments are hard to improve upon.

Beginning late in the \textbf{Report}, and continuing with many ellipses:

\begin{quote}
Is there a right in a government to tax its own funds?

The pretence of this right is deduced from the general right of the legislative power to make all the property of the state contributory to its own exigencies.

\ldots

To tax the funds, is manifestly either to \underline{take}, or to \underline{keep back}, a portion of the principal or interest \underline{stipulated to be paid}.

To do this, on whatever pretext, is \underline{not to do what is expressly promised}; it is not to pay that precise principal, or that precise interest, which has been engaged to be paid; it is, therefore, to violate the promise given to the lender.

But is not the stipulation to the lender, a tacit reservation of the general right of the Legislature to raise contributions on the property of the state?

This cannot be supposed---because it involves two contradictory things; and obligation to do, and a right not to do. An obligation to pay a certain sum, and a right to retain it in the shape of a tax. It is against the rules, both of law and reason, to admit, by implication, in the construction of a contract, a principle which goes in destruction of it.

\ldots

The true definition of public debt is \underline{a property subsisting in the faith of the government}. Its essence is promise. Its definite value depends upon the reliance that the promise will be definitely fulfilled. Can the government rightfully tax its promise? Can it put its faith under contributions? Where or what is the value of the debt, if such a right exists?

\ldots

When a government enters into contract with an individual, it deposes as to the matter of the contract its constitutional authority, and exchanges the character of legislator for that of a moral agent, with the same rights and obligations as an individual. Its promises may be justly considered as excepted out of its power to legislate, unless in aid of them. It is, in theory, impossible to reconcile the two ideas \underline{of a promise which obliges} with \underline{a power to make a law which can vary the effect of it}. This is the great principle that governs the questions, and abridges the general right of the government to lay taxes, excepting out of it a species of property which subsists only in its promise. [----, pp. 159 ff.]
\end{quote}

From an open-loop, or full precommitment point of view, Hamilton's argument makes no sense: there is no difference between paying six percent on untaxed government bonds and twelve percent on bonds taxed at a fifty percent flat rate. Yet it is clear from the quoted passage that Hamilton had time-consistency in mind, that the fact that a bond is a promise issued \underline{today}, while the tax on coupon payments is levied \underline{tomorrow} is at the center of his argument. Our section 2 provides the first analytical context known to us in which Hamilton's argument can be precisely formulated, and within this context it is clear that his argument is correct: If coupon payments are made taxable, in the argument of that section, our time-consistency proof fails, and for exactly the reasons Hamilton set out.

The implications for monetary economies, as developed in section 4, are best seen as a straightforward corollary of Hamilton's argument. In a monetary system the option of monetary expansion offers the opportunity for governments to revise the real ``promise'' intended by bond issues of earlier governments in a way that fundamentally alters the nature of this ``promise.'' For fiscal policy to serve its essential purpose, in dynamic settings, of distributing tax distortions over time in a welfare-maximizing way, monetary policy must be ``tied down'' by, in our example, a commodity standard of a very specific type. Although the form of our optimal standard is obviously specific to the particular model we used to derive it, it seems to us likely that any satisfactory model will exhibit a similar necessity of monetary precommitment of some similar form, and for essentially the same reasons.

We remember Hamilton not only for his financial and economic acuity, but also for his deep distrust of popular democracy and for his close association with a ``creditor class'' who held securities at a time when government policy on taxing the income from such securities was as yet undecided. Historians have, for good reason, been reluctant to accept his arguments on ``faith.'' One of the uses of more specific economic theory is exactly to reduce the need to rely on ``faith'' in evaluating arguments, and in this case, theory has served us effectively. Our formalization of Hamilton's argument has been developed under assumptions that completely remove issues of debtor-creditor conflict from consideration and which express (for analytical purposes only) a kind of \underline{perfect} trust in the motives of government. With these issues swept aside, it is clear that the basic logic does not depend on them in any way, and that the effectiveness, in a modern welfare-economic sense, of fiscal and monetary policy depends in an important way on the ability to make fiscal commitments that are binding in a real sense, and to severely limit the ability of monetary policy to undo such commitments.

\newpage

\section*{Summary of Section 6}

In their concluding section, Lucas and Stokey argue that their theoretical findings were remarkably anticipated by Alexander Hamilton in his 1795 \textit{Second Report on Public Credit}. Hamilton argued that governments should not tax the interest on their own debt, claiming that a bond represents ``a property subsisting in the faith of the government'' and that taxing it would violate the promise to bondholders. From a static perspective, this argument seems puzzling---there should be no difference between paying 6\% on untaxed bonds versus 12\% on bonds taxed at 50\%. However, Lucas and Stokey demonstrate that Hamilton had implicitly understood the time-consistency problem: a bond is a promise made today, while the tax decision comes tomorrow, and this timing distinction is crucial.

The authors show that their Section 2 provides the first rigorous analytical framework in which Hamilton's argument can be precisely formulated and verified. Within their model, if coupon payments were made taxable, the time-consistency of optimal policy would fail for exactly the reasons Hamilton identified. The monetary economy analysis in Section 4 extends this logic: monetary expansion allows governments to effectively revise the real value of past promises through inflation, fundamentally altering the nature of debt contracts. For fiscal policy to optimally distribute tax distortions over time, monetary policy must be constrained---in their model, through a specific commodity standard.

Lucas and Stokey acknowledge Hamilton's historical association with creditor interests and the skepticism this has generated among historians. However, they argue that economic theory allows evaluation of Hamilton's arguments on their logical merits rather than through appeals to faith or suspicion about motives. By constructing a model that eliminates debtor-creditor conflicts and assumes perfect government intentions, they demonstrate that Hamilton's core insight---that effective fiscal and monetary policy requires genuine commitment mechanisms that future governments cannot easily undo---is fundamentally correct and independent of distributional concerns. The time-consistency problem is thus not about trust in government motives, but about the logical structure of intertemporal policy-making itself.

\end{document}
